\documentclass[aspectratio=169]{beamer}
\input{header}
\coursecode{JKSSB}

\title{Current Web, HTML, JavaScript, VBScript \& ASP.NET}
\subtitle{Lesson 63 --- Markup, Scripting and Web History}
\author{JKSSB Computer Awareness}
\date{\today}

\begin{document}
\frame{\titlepage}

\begin{frame}{Where This Lesson Sits}
  \begin{itemize}
    \item A web page can be \key{static} (same content for every visitor) or
      \key{dynamic} (content built or changed at view time).
    \item \key{Markup} describes structure: HTML lays out headings, text, and links.
    \item \key{Scripting} adds behaviour: JavaScript and VBScript run small
      programs inside the page.
    \item \key{Server technology} such as ASP.NET builds pages on the server
      before they reach the browser.
  \end{itemize}
  \begin{block}{The one idea}
    Markup structures the page; scripts make it act; the server can build it first.
  \end{block}
\end{frame}

\begin{frame}{HTML: What It Is}
  \begin{itemize}
    \item \key{HTML} stands for HyperText Markup Language.
    \item It is a \key{markup language}, not a programming language: it tags
      text to say what each part is.
    \item An HTML document is an \key{ASCII-text document}: it can be written
      and read in any plain-text editor.
    \item The browser reads the tags and \key{renders} the page; the tags
      themselves are not shown.
  \end{itemize}
  \begin{block}{Keep the pair straight}
    HTML describes structure; the browser decides appearance from those tags.
  \end{block}
\end{frame}

\begin{frame}{Tags and Attributes}
  \begin{itemize}
    \item A \key{tag} names an element, written in angle brackets such as
      \texttt{<p>} and \texttt{</p>}.
    \item Most elements have an \key{opening tag}, content, and a
      \key{closing tag} with the slash.
    \item An \key{attribute} adds detail inside the opening tag, as
      \texttt{name="value"}, for example \texttt{<body bgcolor="yellow">}.
    \item A \key{valid} tag or attribute is one the language defines; an
      invented name is simply not recognised.
  \end{itemize}
  \begin{exampleblock}{Worked example}
    In \texttt{<font size="4">Hello</font>}, \texttt{font} is the tag and
    \texttt{size="4"} is the attribute.
  \end{exampleblock}
\end{frame}

\begin{frame}{The \texttt{h1} Heading}
  \begin{itemize}
    \item Headings run \key{\texttt{h1}} through \texttt{h6}.
    \item \key{\texttt{h1}} is the \key{largest, top-level heading}; numbers
      grow as headings shrink.
    \item It is written as \texttt{<h1>Title</h1>}.
    \item \muted{Exam depth: \texttt{h1} is biggest; \texttt{h6} is smallest.}
  \end{itemize}
  \begin{alertblock}{Trap alert}
    Do not reverse the order: \key{\texttt{h1}} is largest, not smallest.
  \end{alertblock}
\end{frame}

\begin{frame}{The \texttt{marquee} Tag}
  \begin{itemize}
    \item The \key{\texttt{marquee}} tag makes text \key{scroll} across the page.
    \item It is written as \texttt{<marquee>Welcome</marquee>}.
    \item It is a classic \key{presentation} tag: it moves text but adds no
      structure.
    \item \muted{It belongs to older, browser-specific HTML, not to strict
      modern structure.}
  \end{itemize}
  \begin{block}{One-line memory hook}
    \texttt{marquee} means moving text; \texttt{h1} means the biggest heading.
  \end{block}
\end{frame}

\begin{frame}{DHTML: Dynamic HTML}
  \begin{itemize}
    \item \key{DHTML} stands for Dynamic HyperText Markup Language.
    \item It is \key{not a separate language}: it is HTML plus style sheets
      plus scripts working together.
    \item It lets a page \key{change after loading}, for example text that
      reacts when the mouse passes over it.
    \item Contrast: plain HTML is \key{static}; DHTML pages are
      \key{interactive}.
  \end{itemize}
  \begin{block}{Keep the pair straight}
    HTML lays out a fixed page; DHTML makes that page respond without reloading.
  \end{block}
\end{frame}

\begin{frame}{XML: What It Is}
  \begin{itemize}
    \item \key{XML} stands for eXtensible Markup Language.
    \item It carries and \key{stores data}; it does not lay out a page the way
      HTML does.
    \item Tags are \key{user-defined}: the author invents tags such as
      \texttt{<city>} to label the data.
    \item Contrast: HTML uses a \key{fixed} tag set for display; XML uses an
      \key{open} tag set for data.
  \end{itemize}
  \begin{alertblock}{Trap alert}
    XML stores and transports data; it does not replace HTML for display.
  \end{alertblock}
\end{frame}

\begin{frame}{Java vs JavaScript}
  \begin{center}
    \begin{tabular}{lll}
      \toprule
      \textbf{Point} & \textbf{Java} & \textbf{JavaScript} \\
      \midrule
      Type & Full programming language & Page scripting language \\
      Runs & Compiled to bytecode & Interpreted in the browser \\
      Needs & Java platform & Just a browser \\
      \bottomrule
    \end{tabular}
  \end{center}
  \vspace{0.35cm}
  \begin{alertblock}{Trap alert}
    Similar names, different languages: \key{Java is not JavaScript}.
  \end{alertblock}
\end{frame}

\begin{frame}{JavaScript: History in One Line}
  \begin{itemize}
    \item \key{JavaScript} is a lightweight scripting language for web pages.
    \item It was developed at \key{Netscape} for the Navigator browser in the
      mid-1990s.
    \item It first shipped under other names (Mocha, then LiveScript) before
      the \key{JavaScript} name was settled.
    \item Despite the name, it was \key{not} a cut-down Java; the name was a
      product of that era.
  \end{itemize}
  \begin{block}{Exam anchor}
    JavaScript means Netscape; VBScript means Microsoft.
  \end{block}
\end{frame}

\begin{frame}{JavaScript String Concatenation}
  \begin{itemize}
    \item In JavaScript, the \key{\texttt{+}} operator \key{joins (concatenates)}
      strings.
    \item \texttt{"Good" + "Morning"} gives \texttt{"GoodMorning"}.
    \item Mixing text and a number with \texttt{+} joins as text:
      \texttt{"Total: " + 5} gives \texttt{"Total: 5"}.
    \item \muted{The same symbol adds numbers, but with strings it joins them.}
  \end{itemize}
  \begin{exampleblock}{Worked example}
    \texttt{"Web" + "Page"} evaluates to \texttt{"WebPage"} --- joined, not added.
  \end{exampleblock}
\end{frame}

\begin{frame}{VBScript and \texttt{Dim}}
  \begin{itemize}
    \item \key{VBScript} stands for Visual Basic Scripting Edition, from Microsoft.
    \item It scripts pages and small Windows tasks in a Basic-like style.
    \item Variables are declared with the \key{\texttt{Dim}} statement, short
      for Dimension.
    \item \texttt{Dim city} declares a variable named \texttt{city}.
  \end{itemize}
  \begin{block}{Exam anchor}
    See \texttt{Dim} and think VBScript variable declaration.
  \end{block}
\end{frame}

\begin{frame}{ASP.NET and the \texttt{Session} Object}
  \begin{itemize}
    \item \key{ASP.NET} is Microsoft's framework for building dynamic web pages
      on the \key{server}.
    \item The server builds the page and sends plain HTML to the browser.
    \item The \key{\texttt{Session}} object stores data for \key{one user
      across many pages}, for example a login name.
    \item Contrast: \texttt{Session} is per-user; \texttt{Application} state
      is shared by all users.
  \end{itemize}
  \begin{alertblock}{Trap alert}
    \texttt{Session} keeps one visitor's data; it is not shared across visitors.
  \end{alertblock}
\end{frame}

\begin{frame}{Meta-Search Engines: Dogpile}
  \begin{itemize}
    \item A \key{search engine} keeps its own index; a \key{meta-search engine}
      keeps none and queries other engines.
    \item \key{Dogpile} is the classic exam example of a meta-search engine.
    \item It collects results from several search engines and shows them together.
    \item \muted{Google and Bing are search engines; Dogpile searches the
      search engines.}
  \end{itemize}
  \begin{block}{One-line memory hook}
    Dogpile piles up other engines' results.
  \end{block}
\end{frame}

\begin{frame}{Valid Tags and Attributes: Drill}
  \begin{itemize}
    \item Valid tags include \key{\texttt{h1}}, \texttt{p}, \texttt{marquee},
      \texttt{font}, \texttt{body}, and \texttt{html}.
    \item Valid attributes sit inside the opening tag as
      \key{\texttt{name="value"}}, for example \texttt{size}, \texttt{bgcolor}.
    \item A tag and an attribute are \key{different things}: \texttt{font} is
      a tag; \texttt{size} is its attribute.
    \item Case in the stem does not change the answer: match the \key{name},
      not the capitalisation.
  \end{itemize}
\end{frame}

\begin{frame}{Trap Alert: Facts to Keep Straight}
  \begin{itemize}
    \item HTML is an \key{ASCII-text} document, not a compiled binary.
    \item \key{\texttt{h1}} is largest; \texttt{marquee} scrolls text.
    \item \key{Java} is not JavaScript; \key{XML} carries data, HTML displays it.
    \item \key{\texttt{+}} joins strings in JavaScript; \key{\texttt{Dim}}
      declares in VBScript.
    \item \key{\texttt{Session}} is per-user; \key{Dogpile} is meta-search.
  \end{itemize}
\end{frame}

\begin{frame}{Lesson 63: Recall}
  \begin{itemize}
    \item HTML (HyperText Markup Language) is a markup language stored as an
      ASCII-text document; tags structure, attributes refine.
    \item \texttt{h1} is the largest heading; \texttt{marquee} scrolls text.
    \item DHTML is HTML plus styles plus scripts for pages that change after
      loading; XML is extensible markup for carrying data.
    \item Java is a full compiled language; JavaScript is a Netscape browser
      scripting language, and \texttt{+} concatenates its strings.
    \item VBScript declares variables with \texttt{Dim}; ASP.NET keeps one
      user's data in the \texttt{Session} object.
    \item Dogpile is the classic meta-search engine: it combines other
      engines' results.
  \end{itemize}
\end{frame}
\end{document}
